% Generated draft: requires mathematical and source review.
\documentclass[11pt]{article}
\usepackage[margin=25mm]{geometry}
\usepackage{fontspec}
\setmainfont{Latin Modern Roman}
\usepackage{amsmath,amssymb,mathtools,unicode-math}
\setmathfont{Latin Modern Math}
\usepackage{enumitem,needspace,xurl,hyperref}
\setlength{\parindent}{0pt}
\setlength{\parskip}{5pt}
\setlength{\emergencystretch}{3em}
\newcommand{\N}{\mathbb{N}}
\newcommand{\Z}{\mathbb{Z}}
\newcommand{\Q}{\mathbb{Q}}
\newcommand{\R}{\mathbb{R}}
\newcommand{\C}{\mathbb{C}}
\newcommand{\F}{\mathbb{F}}
\newcommand{\A}{\mathbb{A}}
\newcommand{\PP}{\mathbb{P}}
\newcommand{\E}{\mathbb{E}}
\newcommand{\eps}{\varepsilon}
\newcommand{\dd}{\,\mathrm{d}}
\begin{document}
\section*{UnsolvedMath ID: 701139}
\textbf{Source problem code:} catalogue entry
\subsection{Definitions and mathematical statement}
A group \(G\) is periodic if every element of \(G\) has finite order. Write \([x,y]=x^{-1}y^{-1}xy\), and define the derived subgroup by
\[
[G,G]=\langle [x,y]\mid x,y\in G\rangle.
\]
A \(5\)-group here means a possibly infinite group in which every element has order \(5^n\) for some integer \(n\geq 0\), with \(n\) allowed to depend on the element.

\textbf{Question (V. D. Mazurov, Kourovka Notebook 15.51).} Is it true that, for every periodic group \(G\) satisfying
\[
\forall x,y\in G,\qquad [x,y]^5=1,
\]
the subgroup \([G,G]\) is a \(5\)-group? Equivalently, must
\[
\forall g\in [G,G]\;\exists n\in\mathbb Z_{\geq 0},\qquad g^{5^n}=1
\]
hold? No finiteness assumption on \(G\), and no uniform bound on the orders of elements of \([G,G]\), is imposed.
\subsection{Short English statement}
If every element of a group has finite order and every commutator has order dividing \(5\), must every element of its derived subgroup have order a power of \(5\)?
\subsection{Sources}
\begin{itemize}
\item Local UnsolvedMath registry, supplied record with permanent ID \(701139\), exact source code \texttt{catalogue entry}, upstream revision \texttt{d371ee917ac9647fbb06cd6fc529deec6f3dda82}.
\item \emph{The Kourovka Notebook (Unsolved Problems in Group Theory)}, problem 15.51, proposed by V. D. Mazurov in the 15th issue (2002). The supplied background also identifies the 21st issue edition (2026), \texttt{arXiv:1401.0300v48}. The collection description names E. I. Khukhro and V. D. Mazurov as editors.
\item The supplied registry labels the problem \texttt{open} and records an update timestamp of \texttt{2026-10-16T14:00:00Z}. These are supplied registry metadata, not an independently verified assessment of the literature or mathematical status.
\end{itemize}
\end{document}
